%% ma: L12-B08-0004
%% lop: 12
%% chuong: 2
%% bai: 8
%% dang: 12.8.4
%% loai: DS
%% muc: [TH, TH, VD, VD]
%% dap_an: "SSĐĐ"
%% nguon: "(NBV)-ĐỀ SỐ 2-MỖI NGÀY 1 ĐỀ THI 2025.pdf (D02, MỖI NGÀY 1 ĐỀ - Đề số 2), Phần II Câu 3"
%% kien_thuc: "Bài 8 (tọa độ vectơ, ba điểm thẳng hàng, độ dài vectơ, vectơ cùng phương)"
%% ghi_chu: "Đề gốc có ảnh flycam (không cần). Đáp án gốc S S Đ Đ đúng (kiểm bằng sympy)"
%% trang_thai: chinh-thuc
\begin{ex}
Trong không gian $Oxyz$ với đơn vị trên mỗi trục là $1$ m, một flycam bay với vận tốc có độ lớn và hướng không đổi. Tại thời điểm $t=0$, flycam ở vị trí $A(1;2;3)$ và sau $10$ phút nó ở vị trí $B(21;32;33)$.
\choiceTF
{Flycam không bay qua vị trí $D(5;8;9)$.}
{Vectơ vận tốc của flycam là $\vec v=(20;30;30)$.}
{\True Tốc độ (làm tròn đến hàng phần trăm) của flycam là $0{,}08$ m/s.}
{\True Sau $15$ phút, vị trí của flycam là $C(31;47;48)$.}
\loigiai{Flycam chuyển động thẳng đều từ $A$ đến $B$ trong $10$ phút nên $\overrightarrow{AB}=10\vec v$.
a) $\overrightarrow{AB}=(20;30;30)$, $\overrightarrow{AD}=(4;6;6)=\dfrac15\overrightarrow{AB}$ nên $D$ nằm trên đoạn $AB$, flycam có bay qua $D$. Sai.
b) $\vec v=\dfrac{\overrightarrow{AB}}{10}=(2;3;3)$ (m/phút), không phải $(20;30;30)$. Sai.
c) $|\vec v|=\sqrt{2^2+3^2+3^2}=\sqrt{22}$ m/phút $=\dfrac{\sqrt{22}}{60}\approx0{,}08$ m/s. Đúng.
d) $\overrightarrow{AC}=15\vec v=(30;45;45)\Rightarrow C(31;47;48)$. Đúng.}
\end{ex}
